\documentclass{article}
\usepackage{iclr2026_conference,times}
\usepackage{hyperref}
\usepackage{url}
\usepackage{graphicx}
\usepackage{booktabs}
\usepackage{amsmath}
\usepackage{xcolor}
\usepackage{amsfonts}
\usepackage{amssymb}
\usepackage{amsthm}
\usepackage{enumitem}

% Theorem environments and shared math notation.
\newtheorem{theorem}{Theorem}
\newtheorem{assumption}{Assumption}
\newtheorem{proposition}{Proposition}
\theoremstyle{definition}
\newtheorem{definition}{Definition}
\theoremstyle{plain}
\newcommand{\E}{\mathbb{E}}
\newcommand{\Var}{\mathrm{Var}}
\newcommand{\norm}[1]{\lVert #1 \rVert}

% Tighten float/caption spacing to fit the 9-page limit (formatting only).
\setlength{\textfloatsep}{8pt plus 2pt minus 2pt}
\setlength{\floatsep}{8pt plus 2pt minus 2pt}
\setlength{\intextsep}{8pt plus 2pt minus 2pt}
\setlength{\abovecaptionskip}{4pt}
\setlength{\belowcaptionskip}{2pt}

% -----------------------------------------------------------------------
% Working title — change before submission
% -----------------------------------------------------------------------
\title{The Traits are in the Weights: Estimating Fine-Tuning Lineage in LLMs}

% Authors hidden for double-blind review.
% \iclrfinalcopy  % Uncomment for camera-ready only.

\newcommand{\fix}{\marginpar{FIX}}
\newcommand{\new}{\marginpar{NEW}}

% Placeholder macro for unwritten sections — shows in red, visible in draft
\newcommand{\TODO}[1]{\textcolor{red}{\textbf{[TODO: #1]}}}

% Shorthand for figure path prefix
\newcommand{\figpath}{figures}

\begin{document}
\maketitle

% -----------------------------------------------------------------------
\begin{abstract}
With over a million open models on HuggingFace and a handful of frontier labs constantly updating and releasing models, most models are fine-tuned descendants of only a few ancestors.  However, the true lineage may be undocumented. Model provenance, auditing, and safe deployment all benefit from access to this model lineage, revealing both internal (e.g weights) and external traits (e.g. behavior). We treat fine-tuning as an evolutionary process that leaves signal in weight space, and ask whether we can estimate the model lineage from a set of models. We study lineage at three levels of access: estimation from weight changes alone, the link between those weight changes and behavior, and estimation from behavior alone. We find that the full training tree is recoverable against a known ground-truth tree; that the signal concentrates in a small set of layers; that weight-space proximity predicts behavioral similarity even on a held-out task; and that the same pipeline estimates frontier-model families from outputs with no weight access. Across an encoder--decoder model and decoder-only models from 160M to 8B parameters, we find that lineage estimation is not only possible but is guided by the Atteson margin.

\end{abstract}

% -----------------------------------------------------------------------
\section{Introduction}
\label{sec:intro}




With over a million open models on HuggingFace and a handful of frontier
labs constantly releasing new versions, most models are
fine-tuned descendants of a few ancestors.  Very few of those models document that lineage. Benchmarks can tell us how a model performs on a task, but not how a model originated, how it differs from its sibling models, or how it may
fail. Training lineage can reveal much about model provenance, auditing, and safety~\cite{behavioral2026provenance, nikolic2026model, tensorguard2025}.

We, like other prior works~\citep{horwitz2025mother, yu2025neural, wu2026llmdna}, observe that fine-tuning is a path-dependent, branching process. Each
child model inherits its parent's  updated weights, nudged towards a certain loss and/or data set. This is
precisely an evolutionary process, as described in \cite{yax2025phylolm}, and the evolutionary path leaves signal in weight space. We treat a family of fine-tuned models
as a phylogeny (i.e. an evolutionary tree~\cite{saitou1987neighbor}) and ask whether this \textit{full} training tree can be \emph{estimated} from the model parameters alone.

Consider a firm that forks a base model into finance, IT, and fraud variants, each fine-tuned until
dozens of versions circulate with no recorded lineage (Fig.~\ref{fig:overview}A); a practitioner needs
to know which to build on, where two versions diverged, and where each fails. Phylogenetic algorithms
were built for exactly this inference in biology~\citep{saitou1987neighbor, robinson1981comparison},
and we adapt them to model weights: we estimate a training tree scored against a \emph{known}
ground-truth tree (Fig.~\ref{fig:overview}A), link it to behavior (Fig.~\ref{fig:overview}B), and
compare to behavior-only estimation when weights are unavailable (Fig.~\ref{fig:overview}C).

\begin{figure}[t]
  \begin{center}
    \includegraphics[width=0.95\linewidth]{\figpath/fig1_paper_overview.pdf}
  \end{center}
  \caption{\textbf{Overview.} Fine-tuning leaves a path-dependent signature in weight space
  that we read as a model lineage across three levels of access (decreasing left to right).
  \textbf{(A)}~In controlled trees with known topology, weight distances feed neighbor-joining
  to estimate a phylogeny $\widehat{\mathcal{P}}$, scored against the true tree
  $\mathcal{P}^{*}$ via RF/FN and a random-tree null. \textbf{(B)}~Weight-space proximity
  predicts behavioral similarity, even on a held-out task. \textbf{(C)}~The same pipeline on
  model outputs estimates frontier-model families with no weight access.}
  \label{fig:overview}
\end{figure}




We make three contributions:
\begin{itemize}[leftmargin=1.3em,itemsep=1pt,topsep=2pt,parsep=0pt]
  \item \textbf{A weight-space framework for lineage estimation.} Guided by the geometry of
    fine-tuning updates, our metric and summary-statistic choices estimate the full model lineage; by
    treating layers as `genes' we identify which \emph{carry} the signal (e.g.\ in an encoder--decoder
    model the key projections alone, $4\%$ of parameters, estimate lineage as well as the full tensor).
  \item \textbf{A weight$\to$behavior bridge.} Weight-space proximity predicts behavioral similarity,
    even on a held-out task the models never trained.
  \item \textbf{Behavior-only lineage at frontier scale.} The same pipeline applied to model
    \emph{outputs} estimates lineage for frontier-model families with no weight access.
\end{itemize}
Across all three we span an encoder--decoder model and decoder-only models from 160M to 8B and full
fine-tuning through LoRA, finding that a label-free proxy for the Atteson margin reliably indicates
recovery.

Together, these results show that the training history of a model is legibly
written in its parameters, recoverable as a phylogeny, and predictive of some
behavior.







% -----------------------------------------------------------------------
\section{Related Work}
\label{sec:related}
% ----------------------------------------------------------------------

LLM lineage has become an increasingly studied topic, appearing in three broad forms:
\emph{model-parameter} studies, \emph{weight--behavior} studies that align parameters with
outputs, and \emph{response-output} studies. Across these forms the work is driven by a range of
applications: provenance and intellectual property (IP)
\citep{horwitz2024mother, zeng2024huref, nikolic2026model}, model poisoning and misalignment
\citep{cloud2025subliminal}, inheritance of traits from parent models
\citep{cloud2025subliminal, agents2026same}, and data inquiries and auditing
\citep{zhu2025independence}.

\textbf{Model parameters.} Several works trace parameter changes from fine-tuning or RLHF back to parent models. MoTHer
\citep{horwitz2024mother} recovers family trees from weight similarity via a directed minimum spanning
tree; \citet{yu2025neural} recover parent--child and grandparent--child lineage with both a
learning-free fine-tuning approximation and a learned detector ($>99\%$ accuracy). Beyond direct
lineage, spectral weight ``fingerprints'' identify common ancestry \citep{ghostspec2025}, and Spectral
DeTuning \citep{horwitz2024spectral} recovers base weights from shared LoRA variants. None, however,
recovers a full training tree or evaluates against a known ground-truth tree.

Other works examine \emph{why} lineage carries. Models fine-tuned from a shared initialization stay
within a single low-error basin and are robust to SGD noise
\citep{neyshabur2020,wortsman2022model,frankle2020linear}, so pairwise weight distance is a meaningful
proxy for genealogical distance. Task vectors \citep{ilharco2023editing} --- the displacement a task
induces --- are near-orthogonal across tasks, so distinct lineages contribute largely independent
directions; CKA \citep{kornblith2019similarity} compares models across initializations and
architectures. LoRA \citep{hu2022lora} and intrinsic-dimensionality results
\citep{aghajanyan2021intrinsic} explain why small-rank updates suffice and why rank matters, while
\citet{kirkpatrick2017overcoming} slow learning on weights important to prior tasks.

\textbf{Model parameters $\leftrightarrow$ response output.} A second group links parameter changes to
behavior. \citet{cloud2025subliminal} show child models inherit ``subliminal'' traits even when every
reference to the trait is filtered from training text, and prove a parent's preference is preserved
under a large enough training signal given a shared initialization. REEF \citep{zhang2025reef} compares
internal representations (via CKA) to decide lineage, and gradient-response fingerprints
\citep{tensorguard2025} classify families from how parameters react to probes. These establish that
parameter closeness tracks behavioral similarity, but only pairwise, not at full topology.

\textbf{Response output.} Blackbox studies need only API access or benchmark outputs
\citep{yax2025phylolm, wu2026llmdna}. PhyloLM \citep{yax2025phylolm} builds LLM phylogenies from output
similarity and predicts benchmark performance; LLM DNA \citep{wu2026llmdna} defines a functional
``DNA'' representation carrying lineage; both reconstruct phylogenies over hundreds of models. Other
output-only methods test pairwise provenance from next-token distributions \citep{nikolic2026model} or
output independence \citep{zhu2025independence}. This line is typically unsupervised and prompt-dependent,
with no ground-truth topology to score against.

Our work fills this gap by bridging parameter-based and behavioral lineage while (1) determining where
the training tree can be estimated and (2) evaluating against ground-truth trees. Mirroring
\citet{cloud2025subliminal}, our estimation proposition is: when model-distance signal outweighs training
noise, the lineage can be estimated with high probability.





% -----------------------------------------------------------------------x
\section{Framework}
\label{sec:framework}
% -----------------------------------------------------------------------

In this section, we define our terms and notation and establish the statistical framework underlying our experiments.


Like \citep{wu2026llmdna}, we use phylogenies and LLM ``DNA'' to estimate properties of parent--child training. We assume, like \citep{yu2025neural} but unlike \citep{horwitz2025mother}, that all nodes derive from a single root. Unlike prior work, we estimate partial topology across all nodes, not just parent--child pairs: we form a distance matrix and apply neighbor-joining \citep{saitou1987neighbor} to estimate the tree. We use standard phylogenetics terms throughout: a \emph{phylogeny} is a tree whose leaves are the observed models and whose internal nodes are latent common ancestors, and a \emph{clade} is the set of leaves sharing one such ancestor (equivalently, one branch of the tree). Estimating a clade thus means correctly grouping the models that descend from a common ancestor.


\subsection{Phylogenetic Tree Estimation}
\label{sec:framework:phylo}
% We frame fine-tuning lineage recovery as a phylogenetics problem.
% Pairwise weight distances → NJ tree → RF vs ground truth → permutation p.
% Why NJ > directed MST (MoTHer): statistically principled, known properties.
% Cite: Saitou & Nei 1987 (NJ), Robinson & Foulds 1981 (RF).

To estimate $\widehat{\mathcal{P}}$: (1) choose a set of layers $\mathcal{L}$ and a distance metric $d$;
(2) for each layer $\ell$ compute $D_{i,j}^{(\ell)} = d(L^{(\ell)}_i, L^{(\ell)}_j)$ between models
$i,j$; (3) assemble $\hat{\mathbf{D}}$ as a function of the $D_{i,j}^{(\ell)}$ (e.g.\ the across-layer
average); and (4) estimate $\widehat{\mathcal{P}} = f(\hat{\mathbf{D}})$ with $f$ the NJ method (treating
models as leaves, not internal nodes).

Step (3) lets different layers yield different (or multiple) estimates of $\widehat{\mathcal{P}}$, which we use to evaluate layer importance. Since tree space grows super-exponentially in leaves, exact recovery is unlikely; we use Robinson--Foulds (RF) distance \citep{robinson1981comparison} as a measure of partial recovery --- the number of bipartitions on which two trees disagree. Normalizing by the maximum RF enables comparison across tree sizes: normalized RF near $0$ is near-perfect recovery, near $1$ no better than random.

We assess significance with a permutation test: the per-run $p$-value is the fraction of $1{,}000$
random $n$-leaf trees that match the ground truth at least as well as $\widehat{\mathcal{P}}$, and we
combine the 50 per-run $p$-values with Fisher's method ($-2\sum_i \log p_i \sim \chi^2$). We use both a
random-tree null and a stronger label-permutation null that controls for task-level clustering; full
construction in Appendix~\ref{app:null}.

%% that's t he statistical basis for our method.


\subsection{Distance Metric}
\label{sec:framework:metric}

We use vector-based norms (rather than matrix norms) for portability across architectures, considering
cosine, correlation, $\ell_1$, and $\ell_2$ (the last used in MoTHer~\citep{horwitz2024mother}). We
adopt cosine throughout: it is bounded and scale-invariant across layers, so no single high-magnitude
tensor (e.g.\ the noisy layer norms) dominates, unlike the unbounded $\ell_1$/$\ell_2$. Cosine weight
similarity is also the prevailing choice in recent provenance and fingerprinting work
\citep{zeng2024huref, behavioral2026provenance, zhu2025independence}, keeping our distances comparable
to that literature. For adapter-trained models we compute distances on the \emph{merged} weights---the
base model plus its accumulated low-rank update along the root-to-node path---so every model, whether
fully fine-tuned or LoRA-adapted, is compared on the same full-weight footing.





\subsection{Per Layer Estimated Phylogenies}
\label{sec:framework:layers}
% Each weight tensor → its own NJ tree.
% Distribution of per-layer RF distances is the data; consensus tree is summary.
% More statistically principled than any competitor — flag explicitly.

A methodological contribution is to treat per-layer distance matrices as independent estimates rather
than collapsing to a single all-layer distance, by analogy to biological phylogenies that must cohere
across multiple genes encoding different information \citep{maddison1997gene}. A model with $L$ layers
thus yields $L$ distance matrices, which can each form an estimate of $\widehat{\mathcal{P}}$
(summarized by a consensus tree) or be combined (e.g.\ averaged) into a less noisy $\hat{\mathbf{D}}$.
Beyond a lower-variance estimate, the layer-wise view identifies the most and least important layers by
which change most or least across runs.


\subsection{Guiding Lineage Estimation}
\label{sec:framework:theory}

Estimating lineage poses two questions of differing difficulty. \emph{Ordering} asks only whether
sibling (same-branch) pairs are closer than cousin (cross-branch) pairs; pooling all $O(n^2)$
comparisons per run across $50$ runs, the Mann--Whitney $U$ statistic (equivalently a within-run rank
correlation) detects it even when no single pair is decisive. This is the pairwise-ancestry signal
parent--child methods like MoTHer \citep{horwitz2024mother} and Neural Lineage
\citep{yu2024neurallineage} rely on; we report it in \S\ref{sec:results}. \emph{Lineage} ---
asking neighbor-joining (NJ) to return the full training history --- is harder. NJ is exact under
conditions on tree additivity, which model-parameter distances satisfy only approximately; Atteson's
theorem makes this precise.
\begin{theorem}[\citealp{atteson1999performance}]\label{thm:atteson}
Let $D_{\mathrm{add}}$ be the additive metric of the true topology with internal edge lengths
$\{L_s\}$, and $D$ the observed distance matrix. If
$\;\lVert D-D_{\mathrm{add}}\rVert_\infty < \tfrac12\min_s L_s$, then neighbor-joining applied to $D$
returns the true topology.
\end{theorem}
Writing the bottleneck margin $m=\min_s m_s$ with $m_s=L_s/(2\lVert D-D_{\mathrm{add}}\rVert_\infty)$, an
$m>1$ \emph{guarantees} exact recovery. It is a sufficient, conservative condition --- set by the single
worst edge under a worst-case ($\ell_\infty$) error --- so in practice NJ recovers most splits well
before $m$ reaches $1$; the margin is small for cosine distance when updates are near-orthogonal. Because
$m$ is an \emph{oracle} quantity (forming $D_{\mathrm{add}}$ and the $L_s$ needs the true branch lengths,
which we do not have), we define a label-free proxy computed from $D$ alone, based on the classical
four-point condition for tree metrics \citep{buneman1971recovery}.

\paragraph{Four-point additivity ($A$).} For each quartet the three pairwise-distance sums, sorted
$s_1\ge s_2\ge s_3$, satisfy $s_1=s_2$ under an additive metric \citep{buneman1971recovery}. We average
the resolved fraction over the $Q$ quartets:
\[
  A=\frac{1}{Q}\sum_{q=1}^{Q}\frac{s_2^{(q)}-s_3^{(q)}}{s_1^{(q)}-s_3^{(q)}}\in[0,1] ,
\]
which is $1$ when every quartet is perfectly additive (split resolved) and $0$ at the unresolved,
star-like case, measuring how nearly $D$ satisfies the four-point condition NJ requires; equivalently,
$A=1-\delta$ for the quartet $\delta$-score of \citet{holland2002delta}. We relate $A$ and the oracle
margin to recovery in \S\ref{sec:results}.


% -----------------------------------------------------------------------
\section{Experiments}
\label{sec:experiments}
% -----------------------------------------------------------------------
We structure our experiments around three questions: whether the model lineage can be
estimated from weight changes, whether those changes predict downstream behavior, and
whether family structure is visible from behavioral outputs alone in
frontier models without weight access.

\paragraph{Controlled and real training trees.}
\label{sec:exp:trees}
% Base: Flan-T5-base (250M). Bernoulli branching, 4 tasks, 50 replications,
% 5--8 leaves, known topology by construction.
% 2000 steps, lr=3e-4, batch=8, grad_accum=4, bf16.
% Controlled design enables RF evaluation — all competitors lack this.

To create a controlled experiment, we generate stochastic model training trees beginning from a root model (process detailed in Appendix~\ref{app:details}), producing 50 trees with up to 8 leaf nodes. We randomly assign a task/dataset pair to each edge, and at each split fine-tune a child model from its parent on that pair, giving us the full parameters of each leaf. We then estimate a phylogeny from the \textit{model leaves} and compare it to the known training tree, testing both direct lineage (e.g.\ parent--child relationships) and full lineage estimation.

We vary the root model, task types, and trained layers. We examine an encoder--decoder model
(Flan-T5-base, 250M) and decoder-only models (Pythia-160M, Llama-3.2-1B, and the more frontier-scale
Llama-3.1-8B), fine-tuning on Q\&A, classification, translation, and summarization. For Flan-T5 we run
full fine-tuning and four LoRA \citep{hu2022lora} scopes (key only; q/k/v; all four attention
projections; all projections); the decoder-only models use the same training-tree instructions. We also
validate on two real model lineages from HuggingFace.

\paragraph{Direct lineage.}
Prior methods~\citep{horwitz2024mother, yu2025neural} identify parent--child pairs or correct ordering; ours does too, at the branch level (which models group on a branch). For each run we test whether same-branch leaf pairs are closer in weight space than cross-branch pairs with a one-sided Mann--Whitney $U$ test, summarize the effect with the rank-biserial correlation, and report the within-run point-biserial correlation (Fisher $z$-averaged across runs) between weight distance and the same-branch indicator.

\paragraph{Full topology estimation.}
\label{sec:exp:weights}
Besides direct lineage, we test whether our method estimates the full topology. From \S\ref{sec:framework:theory}, recovery should be possible when the distance matrix is sufficiently additive, for which the label-free four-point additivity $A$ is a practical proxy. For each tree we report RF, FN, the recovery metrics (clade \% recovery and refinement), and \% sig (a within-run permutation test against the label-permutation null), relating them to the oracle Atteson margin and $A$ in \S\ref{sec:results}; a single global $p$-value combines the per-run $p$-values with Fisher's method. The additive Fr\'echet trees also let us identify which layers carry the most inherited signal, via the highest- and lowest-FN layers per architecture.


\paragraph{Predicting held-out behavior.}
\label{sec:exp:behavioral}
We next estimate how strongly whitebox weight distance predicts behavioral similarity on \textit{independent, held-out} probe tasks. We generate a new set of training trees in which translation is held out --- no child model is fine-tuned on translation --- then probe each leaf on translation and embed its text output with a small sentence-transformer (\texttt{all-MiniLM-L6-v2}, $384$-d), repeating each prompt three times per model. We report the random-effects DerSimonian--Laird \citeyear{dersimonianlaird1986} pooled correlation $r$ between weight distance and output-embedding similarity, on the paired per-prompt semantic distance with both quantities $z$-scored within run (placing all models and architectures on one scale). We also probe HellaSwag, which to our knowledge neither Flan nor Llama trained on.

% We then estimate $r$ from the following mixed effects regression:
% \begin{equation}\label{eq:reg}
% d_{\ell_2}(y_{i,b,p}, y_{j,b,p}) \sim r \times d_{cos}\left (\theta_{i,b}, \theta_{j,b} \right )  + \beta_0|b
% \end{equation}
% where $d_{cos}$ is cosine distance, $d_{\ell_2}$ is euclidean distance,  $\theta_{i,b}$ are model parameters from model $i$ and training tree $b$, $y_{i,b,p}$ are the output embeddings from model $i$ from prompt $p$ on training tree $b$, and $\beta_0|b$ is the random intercept per training tree.  The subscripts are to emphasize that we standardize cosine distance for each training and standardize embeddings distance for each training tree and prompt.  The reason for this standardization is to put the model parameters and responses on the same scale.  (CHECK).  We report $r$, its 95\% confidence interval, and $p$-value.

% NOTE: der-simonian..  What is the equation?



\paragraph{Behavior-only lineage (frontier models).}
\label{sec:exp:blackbox}
% GPT, Claude, Gemini, Llama probed via output embeddings across 4 tasks.
% Family structure recoverable when aggregating tasks.
% Single-task confounded by formatting (verbosity, markdown).
% Translation cleanest — language crossing decouples embedding from source content.

Sometimes weights are unavailable and provenance must be inferred from outputs alone. Prior work shows frontier families cluster in behavioral embedding space~\citep{wu2026llmdna, yax2025phylolm}. We extend this by probing 23 frontier and local models (GPT, Claude, Gemini, and open-weight families including Llama and Mistral) across four task types --- summarization, QA, classification, and translation --- with the same neighbor-joining pipeline applied to pairwise sentence-embedding distances between outputs, aggregating the per-prompt trees within and across tasks via the Billera--Holmes--Vogtmann (BHV) Fr\'echet (Sturm) mean in tree space~\citep{billera2001geometry, willis2019confidence} (App.~\ref{app:blackbox}). As in \cite{yax2025phylolm}, this is unsupervised.


% -----------------------------------------------------------------------
\section{Results}
\label{sec:results}
% -----------------------------------------------------------------------

%\subsection{Weights contain both direct ancestry and full lineage}

\paragraph{Direct lineage.}

Direct ancestry is recovered under every condition. The same-vs.-cross-branch ordering test
(Appendix~\ref{app:direct}, Table~\ref{tab:direct}) finds same-branch pairs significantly closer than
cross-branch pairs throughout (one-sided Mann--Whitney $p<10^{-25}$), with rank-biserial correlations
from $0.43$ (Flan LoRA, key only) to $0.96$ (Llama LoRA) and within-run Fisher $z$-averaged correlations
from $-0.60$ to $-0.94$ --- across architectures (Flan-T5, Pythia, Llama), scales (160M--8B), and regimes
(full FT and LoRA). This is the level at which parent--child methods like MoTHer \citep{horwitz2024mother}
and Neural Lineage \citep{yu2025neural, yu2024neurallineage} operate;
we go further to the full topology below.

\paragraph{Full topology estimation.}
\label{sec:results:weighttree}
Our methods also estimate the full tree when the distance matrix is sufficiently additive (a high
Atteson margin; \S\ref{sec:framework:theory}). Aggregated over all $\geq 5$-leaf trees,
Table~\ref{tab:rf_by_leaves} shows Flan-T5-base, Llama-3.2-1B (full FT), and the decoder-only
Pythia-160M control all have low FN ($0.07$--$0.28$) and high clade recovery ($84$--$97\%$); the
low-rank adapters match this on their merged weights (\S\ref{sec:framework:metric}), the Llama-3.2-1B
qkv adapter recovering $94\%$ of clades. Most runs recover the \emph{complete} true topology
(refinement $74$--$95\%$, Table~\ref{tab:rf_by_leaves}). Significance against the harder
label-permutation null is gated by tree size: $0\%$ at $5$ leaves --- where too few within-task leaves
exist to permute for the test to have power --- rising to $57$--$75\%$ at $6$--$8$ leaves, so the
$\geq5$-leaf aggregate ($35$--$50\%$) understates per-run recovery on the trees large enough to test.
Recovery is therefore consistent across architectures, scales, and full-FT vs.\ low-rank training.

% Table 1: RF by leaf count
\begin{table}[t]
\caption{Polytomy-aware tree recovery by architecture and training regime (cosine), aggregated over
  all $\geq 5$-leaf training trees. Truth trees are multifurcating, so any binary reconstruction must
  add resolving edges; we decompose RF into FN (true clades missed --- the real error) and FP
  (binary-forced edges, $=$RF$-$FN). Clade recovery $=1-\mathrm{FN}/|\text{true splits}|$; refinement
  $=$ \% of runs recovering the \emph{entire} true tree (FN$=0$); \% sig $=$ \% of runs beating the
  label-permutation null (the harder within-run test) at $p<0.05$. Pythia-160M (matched to Flan's
  width $d=768$) and Llama-3.2-1B ($d=2048$) isolate architecture and scale from dimensional dilution;
  the two LoRA blocks ($r=8$) show low-rank adapters recover comparably to full FT when compared on
  their merged weights (\S\ref{sec:framework:metric}). Global $p$ is Fisher's combination of the
  per-run permutation $p$-values.}
\label{tab:rf_by_leaves}
\begin{center}
\small
\setlength{\tabcolsep}{4pt}
\begin{tabular}{lrrrrrr}
\toprule
Condition & $n$ & RF (SE) & FN (SE) & Clade rec. & Refinement & \% sig \\
\midrule
\multicolumn{7}{l}{\textit{Full fine-tuning}} \\
Flan-T5-base ($d{=}768$, enc--dec)   & 44 & 1.93 (0.21) & 0.20 (0.08) & 90\% & 84\% & 41\% \\
Pythia-160M ($d{=}768$, dec-only)    & 44 & 1.66 (0.18) & 0.07 (0.05) & 97\% & 95\% & 50\% \\
Llama-3.2-1B ($d{=}2048$, dec-only)  & 40 & 2.05 (0.25) & 0.28 (0.08) & 84\% & 75\% & 35\% \\
\midrule
\multicolumn{7}{l}{\textit{LoRA ($r{=}8$, merged weights)}} \\
Flan-T5-base qkv (enc--dec)          & 43 & 2.12 (0.26) & 0.30 (0.09) & 82\% & 74\% & 37\% \\
Llama-3.2-1B qkv (dec-only)          & 44 & 1.80 (0.20) & 0.14 (0.07) & 94\% & 91\% & 48\% \\
\bottomrule
\end{tabular}
\end{center}
\end{table}

Fig.~\ref{fig:behavioral_vs_whitebox} plots FN vs.\ tree size for Flan-T5 and Llama-3.2-1B, with
translation held out so the behavioral series (red) is a fair probe, alongside whitebox weight recovery
(blue) and the random null (gray). Whitebox FN stays low at every size and below both the behavioral and
null series: weight access, not behavior, is what enables per-run topology recovery.

% Figure 2: behavioral vs whitebox vs null
\begin{figure}[t]
  \begin{center}
    \includegraphics[width=0.82\linewidth]{\figpath/fig2_fn_multi.png}
  \end{center}
  \caption{FN (true clades missed) by tree size for whitebox weight distances (cosine, blue),
    held-out behavioral distances (translation probe, semantic, red), and the random null (gray;
    bands mean $\pm$ 2\,SE), for an encoder--decoder (Flan-T5-base) and a decoder (Llama-3.2-1B,
    full FT), both with translation held out of training. FN (vs.\ total RF) removes the
    binary-forced FP the multifurcating truth induces equally for signal and null. In both
    architectures whitebox FN is $\approx 0$ at every tree size, while the behavioral probe tracks
    the null: behavioral similarity is a pairwise lineage gradient, not a recoverable topology.
    Weight access is what enables per-run tree recovery.}
  \label{fig:behavioral_vs_whitebox}
\end{figure}

Treating layers as independent estimators (Appendix~\ref{app:layertrace},
Fig.~\ref{fig:layertrace_grid}), attention key projections are the dominant carrier in both
architectures: in Flan-T5 the key projection alone (fewer than $4\%$ of parameters) recovers lineage at
least as well as the full tensor, and the decoder-only Llama mirrors this key-matrix carrier --- without
cross-attention, so it is not an encoder--decoder artifact --- though more diffusely. Layer norms carry
the weakest signal.

Figure~\ref{fig:coherence_recovery} explores clade recovery across LoRA settings, datasets, layers, and
models. The oracle Atteson margin (left) governs recovery reliably; because it is a sufficient condition
applied stochastically, recovery can be high below margin $1$ and short of $100\%$ above it. Architecture
and layer differences remain visible: single self-attention q/k/v matrices from Llama LoRA recover only
$15$--$40\%$ of clades, versus higher recovery for Flan cross-attention q/k/v, and full-FT layers
($>80\%$) carry more than LoRA adapters. The right panel shows the label-free proxy $A$ predicts clade
recovery adequately, making it a practical tool.

\begin{figure}[t]
  \begin{center}
    \includegraphics[width=0.78\linewidth]{\figpath/fig4_coherence_atteson.png}
  \end{center}
  \caption{\textbf{What gates topology recovery.} \textbf{(A)} Clade recovery against the oracle
    bottleneck Atteson margin (min.\ fitted internal edge length over twice the non-additivity
    error); NJ is guaranteed to recover a split once the margin exceeds $1$ (dashed line).
    \textbf{(B)} The same recovery against the label-free quartet additivity proxy $A$, which needs
    no ground-truth tree. Small points are individual Llama-1B layers or q/k/v adapter matrices;
    large symbols are whole-run summaries, Flan LayerTrace subsets, organic HF validation, and
    layer-family centroids.}
  \label{fig:coherence_recovery}
\end{figure}

Table~\ref{tab:variants} summarizes signal and recovery across regimes: both the ordering signal and
topology recovery are stable from the narrowest LoRA scope (key only) through full fine-tuning. The
within-run ordering correlation strengthens modestly with more trained parameters, while topology
recovery stays high throughout (clade recovery $82$--$96\%$, refinement $72$--$94\%$). Notably,
Llama-3.2-1B LoRA adapters recover as well as or better than Flan full FT --- low-rank adaptation does
not cost lineage recoverability. The stronger within-run ordering correlation for the Llama adapters
than for Llama full FT reflects the \textsc{trained\_only} restriction: the adapter distance is read
from the q/k/v updates alone, whereas the full-FT distance averages over all tensors, including the
noisier layer norms and embeddings.

% Table 2: training-regime variants (Flan-T5 + Llama-3.2-1B)
\begin{table}[t]
\caption{Lineage signal and recovery by training regime (cosine; \textsc{trained\_only=1} for LoRA,
  Llama LoRA distances on merged weights). \emph{Ordering:} Fisher-$z$ $r$ is the within-run correlation
  between cosine distance and the same-branch indicator (SD over runs, $\geq 4$ leaves).
  \emph{Topology:} RF and FN (true clades missed) are polytomy-aware NJ means (SE in parentheses);
  clade recovery $=1-\mathrm{FN}/|\text{true splits}|$, refinement $=$ \% runs with $\mathrm{FN}=0$.
  Every condition has global Fisher-combined $p<10^{-5}$. What gates recovery is analysed in
  Fig.~\ref{fig:coherence_recovery}.}
\label{tab:variants}
\begin{center}
\small
\setlength{\tabcolsep}{4pt}
\begin{tabular}{lrrrrrr}
\toprule
Variant & $n$ & Fisher-$z$ $r$ (SD) & RF (SE) & FN (SE) & Clade rec. & Refinement \\
\midrule
\multicolumn{7}{l}{\textbf{Encoder--decoder: Flan-T5-base} ($d=768$)} \\
\multicolumn{7}{l}{\quad\textit{LoRA adapter scope}} \\
\quad k\_only          & 47 & $-0.597$ (0.33) & 2.11 (0.25) & 0.34 (0.09) & 82\% & 72\% \\
\quad qkv              & 46 & $-0.670$ (0.25) & 1.98 (0.25) & 0.28 (0.08) & 83\% & 76\% \\
\quad full\_attn       & 42 & $-0.671$ (0.25) & 1.83 (0.25) & 0.21 (0.07) & 86\% & 81\% \\
\quad full\_ft\_approx & 40 & $-0.687$ (0.24) & 1.85 (0.25) & 0.23 (0.08) & 88\% & 82\% \\
\multicolumn{7}{l}{\quad\textit{Full fine-tuning}} \\
\quad Tree sample      & 47 & $-0.802$ (0.23) & 1.81 (0.21) & 0.19 (0.07) & 90\% & 85\% \\
\midrule
\multicolumn{7}{l}{\textbf{Decoder-only: Llama-3.2-1B} ($d=2048$)} \\
\quad LoRA qkv ($r{=}8$)  & 47 & $-0.943$ (0.13) & 1.68 (0.20) & 0.13 (0.07) & 94\% & 91\% \\
\quad LoRA qkv ($r{=}64$) & 47 & $-0.960$ (0.11) & 1.60 (0.19) & 0.09 (0.05) & 96\% & 94\% \\
\quad full FT             & 44 & $-0.690$ (0.33) & 1.86 (0.24) & 0.25 (0.07) & 85\% & 77\% \\
\bottomrule
\end{tabular}
\end{center}
\end{table}

Finally, we move to organic model families on HuggingFace. Figure~\ref{fig:real-trees} compares our
reconstruction against the documented lineage for ten public Llama-2-7B checkpoints whose relationships
are recorded on their model cards. Using cosine distances over the common tensors and the
Fr\'echet (Sturm) mean, we recover the single informative split --- the seven chat-derived models versus
the base/Nous-Hermes/Vicuna group --- exactly ($100\%$ clade recovery, $95.2\%$ Fr\'echet support). The
pipeline we validate against synthetic ground truth thus reconstructs ancestry in an uncontrolled
setting with heterogeneous, sharded checkpoints and undocumented intermediate training. (An organic
Flan-T5 depth-3 branch recovers its informative clade similarly.)

\begin{figure}[t]
  \begin{center}
    \includegraphics[width=0.72\linewidth]{\figpath/hf_llama2_trees_vs_truth.png}
  \end{center}
  \caption{\textbf{Recovery on an organic HuggingFace model family.} Reconstructed phylogeny (left;
    cosine Fr\'echet mean of weight distances over tensors common to all ten checkpoints)
    against the documented lineage (right) for a real Llama-2-7B family. Nodes are coloured by side
    of the one informative split: \texttt{Llama-2-7b-chat} and its six descendants (blue) versus
    base, Nous-Hermes, and Vicuna (orange). The reconstruction treats all models as leaves but
    recovers that split exactly; remaining edges are binary-forced resolutions of the polytomous
    truth (\S\ref{sec:framework:phylo}), not lineage errors.}
  \label{fig:real-trees}
\end{figure}


\paragraph{Predicting held-out behavior.}
\label{sec:results:behavioral}

Estimating lineage from weights does not by itself imply that parameter-space proximity tracks
behavioral proximity. Inspired by \citet{cloud2025subliminal}'s subliminal-learning experiment, we test
whether models closer in weight space produce more similar outputs on a held-out task, using translation
(on a tree that never trained on it) and HellaSwag (untrained by either model).
Table~\ref{tab:behavior_holdout} reports a weak but significant negative correlation for HellaSwag in
both Llama and Flan, and for translation in Llama (the sign reverses for Flan translation) --- supporting
\citet{cloud2025subliminal} that unexpected traits can be passed to child models. Full regression
coefficients are in Appendix~\ref{app:regression}, and the degenerate probe behavior behind the Flan
reversal is detailed in Appendix~\ref{app:translation_probe}.





\begin{table}[t]
\caption{Weight$\to$behavior bridge by model and probe (full fine-tuning). Random-effects
(DerSimonian--Laird) pooled correlation $r$ between weight distance and behavioral similarity
(paired per-prompt semantic distance, both $z$-scored within run, so $r$ is comparable across rows).
``trained'' $=$ probe task in the training mix; ``held-out'' $=$ no node trained on it (HellaSwag is
held out for every model).}
\label{tab:behavior_holdout}
\centering
\small
\begin{tabular}{llccc}
\toprule
Model & Probe (condition) & DL $r$ & 95\% CI & $p$ \\
\midrule
Llama-3.2-1B (full FT) & Translation (trained)  & $-0.294$ & $[-0.41, -0.17]$ & $<10^{-5}$ \\
                       & Translation (held-out) & $-0.333$ & $[-0.43, -0.23]$ & $<10^{-5}$ \\
                       & HellaSwag (held-out)   & $-0.204$ & $[-0.32, -0.09]$ & $9\times10^{-4}$ \\
\midrule
Flan-T5-base (full FT) & Translation (trained)  & $-0.260$ & $[-0.35, -0.16]$ & $<10^{-5}$ \\
                       & Translation (held-out) & $+0.113$ & $[+0.02, +0.21]$ & $0.024$ \\
                       & HellaSwag (held-out)   & $-0.292$ & $[-0.37, -0.21]$ & $<10^{-5}$ \\
\bottomrule
\end{tabular}
\end{table}


\paragraph{Behavior-only lineage (frontier models).}
\label{sec:results:blackbox}

Behavior-alone clade recovery (Fig.~\ref{fig:behavioral_vs_whitebox}) beats random but is far less
reliable than weight-based recovery. With only inference access, we estimate lineage from 23 models as a
Fr\'echet tree averaged across four prompt tasks (Fig.~\ref{fig:blackbox_hero}): the tightest cluster is
the GPT branch (though Sonnet infiltrates it), and open models appear more similar to each other than the
frontier models, except Gemini. The estimate is not robust across tasks --- pairwise normalized RF among
the four per-task trees is $0.65$--$0.95$ (mean $\approx 0.84$), so no single task reproduces the
aggregate (Appendix~\ref{app:blackbox}, Fig.~\ref{fig:frechet_grid}) --- and, lacking a ground-truth
baseline, the recovered lineage looks reasonable but is highly prompt-dependent.


% Hero figure: cross-task Frechet phylogeny + BHV tree-space MDS
\begin{figure}[t]
  \begin{center}
    \includegraphics[width=0.72\linewidth]{\figpath/hero_phylogeny_treespace.png}
  \end{center}
  \caption{\textbf{A behavioral phylogeny of frontier and open models, with its support in tree space.}
    \textbf{(a)} BHV Fr\'echet (Sturm) mean of $n=245$ per-prompt neighbor-joining trees (output-embedding
    distances over the 245/315 prompts all 23 models answered, four task families), coloured by
    developer and rendered as a cladogram (internal edges $\sim$3\% of total length). Claude Sonnet and
    Haiku form the most stable clade; the GPT family and Llama pair are also recovered. \textbf{(b)} The
    17 per-dataset trees embedded in BHV tree space (classical MDS of geodesic distances). Trees group
    by task---summarization tightest, translation the outlier---and the Fr\'echet mean ($\star$) lies at
    the centre.}
  \label{fig:blackbox_hero}
\end{figure}




% -----------------------------------------------------------------------
\section{Discussion}
\label{sec:discussion}
% -----------------------------------------------------------------------

We show that the model lineage can be estimated from parameters alone; the best layers --- the
$k$-matrices, only $4\%$ of parameters --- estimate trees as well ($80$--$90{+}\%$) as the full weight
set, this holds across encoder--decoder and decoder models and survives LoRA at low rank ($r=8$),
and every controlled condition is significant. The weights also carry less obvious traits: on held-out
tasks, models closer in weight space produce more similar outputs (pooled DL $r\approx-0.3$ to $-0.2$)
--- modest but direct evidence of inheritance, stronger under full fine-tuning than LoRA. Behavior-only
lineage is informative but prompt-dependent and less robust than parameter estimates, which need no
prompts.

Our results further corroborate that attention layers are critical to LLM learning, here for ancestry,
with cross-attention a stronger lineage carrier than self-attention (to be validated on more
architectures). Update geometry matters --- near-orthogonal updates are hard to recover, one reason
cosine distance is central --- and the work foregrounds full-access vs.\ inference-only auditing: full
access gives the best recovery, yet weight changes track behavioral changes even on held-out tasks.

Limitations: parameter-based estimation currently requires an identical architecture (so Llama cannot be
compared directly to Mistral), and our controlled tree-generation is simple, which we partly mitigate
with the real-world example. Future work includes relating models across architectures (e.g.\ via
CKA~\citep{kornblith2019similarity}) and probing the mechanism of inheritance. We hope this work
illuminates the need to consider models holistically, not as a series of benchmarks.


% -----------------------------------------------------------------------
% Author acknowledgments omitted for double-blind review; add at camera-ready.

\subsubsection*{Use of Large Language Models}
In accordance with the ICLR policy on LLM usage, we disclose that LLM-based coding and writing
assistants were used substantially during this project. Their uses included writing and debugging
code throughout the experimental pipeline---data processing, weight-distance computation, tree
reconstruction and scoring, statistical analysis, and figure and table generation---as well as editing
the manuscript (condensing prose, standardizing terminology, and formatting). The research questions,
experimental design, and scientific claims are the authors' own, and the authors reviewed and verified
all code, numerical results, figures, and citations, taking full responsibility for the content.

% -----------------------------------------------------------------------
\bibliography{references}
\bibliographystyle{iclr2026_conference}


% -----------------------------------------------------------------------
\appendix
% -----------------------------------------------------------------------

\section{Experimental Details}
\label{app:details}

\subsection*{Tree Generation}
\label{app:tree-gen}

Training trees are generated via a Poisson branching process as follows:

\begin{enumerate}
    \item The root node (pretrained Flan-T5-base) draws $k \sim \mathrm{Poisson}(\lambda = 1.5)$ children, with a minimum of 1.
    \item Each subsequent node draws $k \sim \mathrm{Poisson}(\lambda = 1.5)$ children. Nodes that draw $k = 0$ naturally become leaves --- no separate pruning step is needed.
    \item Breadth First Search expansion continues until either a node budget ($n_\mathrm{nodes} = 14$ training nodes, not counting the root) is exhausted or a maximum depth ($n_\mathrm{layers} = 4$) is reached. The budget cap creates the truncation; any node alive when the budget runs out becomes a leaf.
    \item Each non-root node is assigned a task and dataset sampled without replacement from 4 task types (summarization, classification, QA, translation) and 36 datasets (9 summarization, 10 classification, 8 QA, 9 translation). The topology (parent--child structure) is fixed by a topology seed; dataset assignments vary across the 50 replications.
    \item Runs with fewer than 4 leaves are excluded from analysis, as trees with 3 or fewer leaves have trivial topology with no meaningful RF distance.
\end{enumerate}

This procedure yields trees with 4--8 observed leaves across the 50 replications, with variable depth and branching structure. The ground-truth topology is fully known by construction, enabling exact RF evaluation.

\subsection*{Training Configuration}
Base model: Flan-T5-base (250M parameters). Each node trained for 2,000 steps, learning rate $3 \times 10^{-4}$, batch size 8, gradient accumulation 4 (effective batch 32), warmup 200 steps, weight decay 0.01, bf16 precision. Up to 10,000 training examples and 1,000 evaluation examples per dataset, fixed dataset seed across all nodes to avoid data variation as a confound.

\section{Layer Variant Definitions}
\label{app:layers}
Two families of variants appear in the paper. \emph{LayerTrace subsets} (Table~\ref{tab:layer_subsets})
restrict the full-weight distance to a chosen set of tensors, testing which layers carry the lineage
signal; they are defined on Flan-T5-base (12 encoder + 12 decoder blocks). \emph{LoRA variants}
(Table~\ref{tab:lora_variants}) instead change which modules are fine-tuned, all at rank $r=8$,
$\alpha=16$, dropout $0.05$, with adapters merged into the base weights before analysis.

\begin{table}[h]
\caption{LayerTrace layer subsets (Flan-T5-base). ``SA.k'' is the self-attention key projection;
``pre-attention LayerNorm'' is the \texttt{layer.0.layer\_norm} preceding self-attention. Percentages
are of the full trained-tensor set. \texttt{low\_signal} is a negative control (the noisiest LayerNorms).}
\label{tab:layer_subsets}
\centering
\begin{tabular}{lll r}
\toprule
Variant & Tensors included & Count & \% \\
\midrule
\texttt{full}        & all trained tensors                                   & ---  & 100 \\
\texttt{high\_signal}& SA.k $+$ pre-attention LayerNorm (encoder $+$ decoder) & 48   & $\sim$17 \\
\texttt{sak}         & SA.k only (encoder $+$ decoder)                       & 24   & $\sim$8.5 \\
\texttt{encoder\_h}  & encoder SA.k $+$ encoder pre-attention LayerNorm      & 24   & $\sim$8.5 \\
\texttt{enc\_sak}    & encoder SA.k only                                     & 12   & $\sim$4 \\
\texttt{low\_signal} & post-FFN / final-block LayerNorms (control)           & 36   & --- \\
\bottomrule
\end{tabular}
\end{table}

\begin{table}[h]
\caption{LoRA target-module variants (Flan-T5-base, $r=8$, $\alpha=16$, dropout $0.05$). Modules:
$q,k,v,o$ = attention query/key/value/output projections; $w_i,w_o$ = FFN input/output projections.}
\label{tab:lora_variants}
\centering
\begin{tabular}{ll r r}
\toprule
Variant & Target modules & Trainable params & \% of model \\
\midrule
\texttt{lora\_k\_only}        & $k$                       & $\sim$98K  & 0.04 \\
\texttt{lora\_qkv}            & $q,k,v$                   & $\sim$294K & 0.12 \\
\texttt{lora\_full\_attn}     & $q,k,v,o$                 & $\sim$393K & 0.16 \\
\texttt{lora\_full\_ft\_approx}& $q,k,v,o,w_i,w_o$        & $\sim$2.1M & $\sim$0.87 \\
\bottomrule
\end{tabular}
\end{table}

\subsection*{LayerTrace: Per-Layer Recovery}
\label{app:layertrace}
Treating each layer as an independent estimator, Fig.~\ref{fig:layertrace_grid} reports the per-layer
FN-rate across architectures. The attention key matrices recur once per block and miss the fewest true
clades (10\% vs.\ 14\% elsewhere in Flan; 13\% vs.\ 20\% in Llama), while \textsc{LayerNorm} is the
periodic maximum. In Flan the encoder \textsc{SelfAttention.k} and encoder--decoder cross-attention
keys are lowest; the decoder-only model recovers the same key-matrix carrier without any
cross-attention. Across blocks every block contributes, with slight evidence that earlier blocks carry
somewhat more signal.

\begin{figure}[h]
  \begin{center}
    \includegraphics[width=\linewidth]{\figpath/fig3_layertrace_grid.png}
  \end{center}
  \caption{Per-layer FN-rate (\% of true clades missed; mean $\pm$ 2\,SE, cosine, full
    fine-tuning) for an encoder--decoder (Flan-T5-base, $d{=}768$, $n=44$) and a decoder-only model
    (Llama-3.2-1B, $d{=}2048$, $n=47$). The same layers carry lineage in both: the attention
    \emph{key} matrices recur once per block and miss the fewest true clades (10\% vs.\ 14\%
    elsewhere in Flan; 13\% vs.\ 20\% in Llama), while \textsc{LayerNorm} is the periodic maximum
    (noise floor). The decoder-only model recovers the same key-matrix carrier without any
    cross-attention, so the localization is not an encoder--decoder artifact.}
  \label{fig:layertrace_grid}
\end{figure}

\section{Direct-Lineage Ordering Statistics}
\label{app:direct}
Table~\ref{tab:direct} gives the full same-vs.-cross-branch ordering statistics summarized in
\S\ref{sec:results}, for each training regime.

\begin{table}[h]
\caption{Direct-lineage ordering signal by training regime (cosine weight distance, same-vs.\
  cross-branch at the first split; \textsc{trained\_only=1} for LoRA groups). $U$-test is a
  one-sided Mann--Whitney on same-branch ($n_{\text{same}}$) vs.\ cross-branch ($n_{\text{cross}}$)
  pairwise distances pooled over runs; rank-biserial is its effect size (larger $=$ siblings
  closer). Within-run $r$ is the point-biserial correlation between distance and the same-branch
  indicator, Fisher $z$-averaged over the $N$ runs with within-run variation ($\geq 4$ leaves).
  All Llama LoRA distances are computed on merged weights (\S\ref{sec:framework:metric}), as throughout.}
\label{tab:direct}
\begin{center}
\begin{tabular}{llrrr}
\toprule
Model (regime) & $n_{\text{same}}/n_{\text{cross}}$ & MW $p$ & rank-bis. & within-run $r$ ($N$) \\
\midrule
Flan-T5-base (full FT)      & $386/361$ & $<10^{-5}$  & $0.61$ & $-0.81$ (35) \\
Flan-T5-base (LoRA k\_only) & $416/377$ & $<10^{-5}$  & $0.43$ & $-0.60$ (37) \\
Flan-T5-base (LoRA qkv)     & $409/369$ & $<10^{-5}$  & $0.53$ & $-0.67$ (36) \\
Pythia-160M (full FT)       & $401/377$ & $<10^{-5}$  & $0.77$ & $-0.86$ (37) \\
Llama-3.2-1B (full FT)      & $406/339$ & $<10^{-5}$  & $0.54$ & $-0.69$ (35) \\
Llama-3.2-1B (LoRA qkv)     & $416/377$ & $<10^{-5}$ & $0.96$ & $-0.94$ (37) \\
Llama-3.1-8B (LoRA qkv)     & $414/374$ & $<10^{-5}$ & $0.96$ & $-0.93$ (37) \\
\bottomrule
\end{tabular}
\end{center}
\end{table}

\section{Null Distribution Construction}
\label{app:null}

We assess topology recovery against two permutation nulls, each computed per run (per tree) and then
combined across the $50$ runs of a condition.

\paragraph{Random-tree null.} For a run with $n$ leaves we draw $1{,}000$ trees uniformly from the space
of binary topologies on the same $n$ labels and score each against the ground-truth training tree with
the same polytomy-aware recovery metric used for the reconstruction (Robinson--Foulds / clade
recovery; \S\ref{sec:framework:phylo}). The per-run $p$-value is the fraction of null trees whose
distance to the truth is at most the observed distance---the probability that a random tree matches the
truth as well as the reconstruction does. This null asks only whether the reconstruction beats chance.

\paragraph{Label-permutation null.} The random-tree null can be passed by a reconstruction that merely
recovers coarse task-level structure (all summarization leaves in one clade, all classification leaves
in another) without resolving lineage \emph{within} a task block. To control for this we use a stronger
null that destroys only the within-block lineage information: holding the task-block structure fixed, we
randomly permute the lineage labels among leaves that share a task, re-score against the truth, and take
the fraction of permutations at least as good as observed as the $p$-value. ``\% significant vs.\
label'' is the share of runs whose reconstruction beats this permutation at $\alpha=0.05$; it is the
criterion for recovering within-task lineage rather than task identity alone, and is the harder test
reported throughout.

\paragraph{Combining across runs.} Treating the $50$ runs as independent trees, we pool their per-run
$p$-values with Fisher's method, $-2\sum_{i}\ln p_i \sim \chi^2_{2k}$ under the global null of no lineage
information, giving a single experiment-level $p$-value per condition.

\section{Blackbox Experiment Details}
\label{app:blackbox}

\paragraph{Model collection.} The blackbox phylogeny is built over a collection of frontier and open
models spanning the major developer families: OpenAI (GPT-3.5, GPT-4o / 4o-mini, GPT-4.1, o-series),
Anthropic (Claude 3 / 3.5 / 4 Sonnet and Haiku), Google (Gemini 1.5 / 2.0 / 2.5, Gemma), Meta (Llama 2
and 3), Mistral, DeepSeek, and Microsoft Phi. Proprietary models are queried through their providers'
APIs; open-weight models are run locally from HuggingFace checkpoints. The analyses in
\S\ref{sec:results} use the $23$ models that returned valid responses on all four task
families, so that every per-task tree is built over a common leaf set.

\paragraph{Prompts and decoding.} Each model answers a shared bank of $315$ prompts drawn from four task
families---summarization, question answering, classification, and translation. Prompts are issued with
identical text to every model and decoded greedily (temperature $0$, or the provider's nearest
deterministic setting) so that differences in output reflect the model, not sampling noise. Of the
$315$ prompts, $245$ were answered by all $23$ common models and are used for the cross-task aggregate.

\paragraph{Embedding and distances.} Each response is embedded with the
\texttt{all-MiniLM-L6-v2} sentence-transformer ($384$-dimensional, mean-pooled), giving a vector per
(model, prompt). For each prompt we form the pairwise cosine-distance matrix over models and build a
neighbor-joining tree, yielding one tree per prompt. This per-prompt construction decouples the
embedding from any single source text and is what makes translation a clean probe (the language change
separates output style from source content).

\paragraph{Aggregation.} Per-prompt trees are summarized with the Billera--Holmes--Vogtmann (BHV)
Fr\'echet (Sturm) mean in tree space~\citep{billera2001geometry, willis2019confidence}, both within
each task family and across all four (Fig.~\ref{fig:blackbox_hero}, Fig.~\ref{fig:frechet_grid}). We report the aggregate as a cladogram
because the Fr\'echet mean's internal edges are short relative to total edge length, so topology is more
informative than scaled branch length.

\begin{figure}[h]
  \begin{center}
    \includegraphics[width=\linewidth]{\figpath/frechet_grid_tasks.png}
  \end{center}
  \caption{\textbf{Per-task behavioral phylogenies (BHV Fr\'echet means).} The cross-task aggregate
    (top-left) and the four per-task Fr\'echet trees (summarization, QA, classification, translation),
    each over the 23 models common to all tasks and rendered as cladograms. Topologies differ
    substantially across tasks (pairwise normalized Robinson--Foulds $0.65$--$0.95$), which motivates
    the multi-task aggregation reported in \S\ref{sec:results}.}
  \label{fig:frechet_grid}
\end{figure}

\section{Behavioral Regression: Full Results}
\label{app:regression}

All models regress behavioral output similarity (sentence-embedding cosine, translation probe
unless noted) on cosine weight distance with run fixed effects and SEs clustered on run. We report
three coefficients per condition (Table~\ref{tab:regression_betas}):
\begin{itemize}
  \item \textbf{Raw $\beta$} --- the unscaled slope. Significant and negative everywhere, but
    \emph{not} comparable across conditions: LoRA groups have far smaller absolute cosine distances,
    so their raw $\beta$ is inflated by orders of magnitude (e.g.\ $-52$ for full\_ft\_approx vs.\
    $-0.19$ for k\_only) without implying a larger effect.
  \item \textbf{Semi-standardized $\beta$} --- predictor $z$-scored within run; reads as ``$\Delta$
    output similarity per 1 within-run SD of weight distance.'' Interpretable within a condition,
    but still not cross-comparable because conditions differ in how much their \emph{outputs} vary
    (k\_only translations are near-identical, so its semi-standardized $\beta$ is tiny despite the
    strongest correlation).
  \item \textbf{Standardized $\beta$ (within-run)} --- both sides $z$-scored within run. This is the
    only cross-condition-comparable effect size, and it equals the within-run correlation. It is the
    number reported in the main text.
\end{itemize}

\begin{table}[h]
\caption{Behavioral regression coefficients by condition (run FE, cluster-robust SE on run; $p$ for
  the standardized $\beta$). The standardized $\beta$ (within-run) is the comparable effect size and
  equals the within-run correlation; raw and semi-standardized $\beta$ are reported for transparency.
  \texttt{no\_translation\_v2} is a qualitative null (Appendix~\ref{app:translation_probe}) and is
  omitted. HellaSwag is a held-out task for Llama (never trained).}
\label{tab:regression_betas}
\begin{center}
\begin{tabular}{lrrrr}
\toprule
Condition & Raw $\beta$ & Semi-std $\beta$ & Std $\beta$ (within-run) & $p$ \\
\midrule
\multicolumn{5}{l}{\textit{Flan-T5-base}} \\
lora\_k\_only          & $-0.19$  & $-0.0004$ & $-0.41$ & $5\!\times\!10^{-29}$ \\
lora\_qkv              & $-16.6$  & $-0.024$  & $-0.18$ & $2\!\times\!10^{-6}$ \\
lora\_full\_attn       & $-32.0$  & $-0.027$  & $-0.12$ & $0.008$ \\
lora\_full\_ft\_approx & $-52.2$  & $-0.031$  & $-0.14$ & $0.007$ \\
Tree sample (full FT)  & $-7.40$  & $-0.019$  & $-0.13$ & $0.002$ \\
\midrule
\multicolumn{5}{l}{\textit{Llama-3.1-8B (decoder-only; merged weights)}} \\
k\_only (translation, seen)   & $+0.0002$ & $-0.0001$ & $-0.01$ & $0.76$ \\
k\_only (HellaSwag, held-out) & $-0.012$  & $-0.0028$ & $-0.16$ & $0.003$ \\
\bottomrule
\end{tabular}
\end{center}
\end{table}

The coefficients above use the per-run \emph{centroid} semantic DV (one embedding per node). The
held-out bridge in the main text (Table~\ref{tab:behavior_holdout}) instead uses the \emph{paired}
per-prompt semantic distance, which removes the shared prompt-content confound and strengthens the
association in every group; we report the paired DV in the main text and keep the centroid DV here for
continuity with the per-condition effect sizes. The Llama rows here are the 8B decoder, included as an
additional scale point; the headline held-out bridge in Table~\ref{tab:behavior_holdout} rests on the
full-FT Llama-3.2-1B and Flan-T5 models. The two agree where they overlap --- held-out HellaSwag is
negative and significant for the 8B decoder here ($-0.16$) as for the 1B in the main text ($-0.20$).

\paragraph{Pooling estimator.} Because each run has its own tree and hence its own weight-space scale, a
raw pooled slope is not meaningful across runs; we therefore standardize within run (equivalently, treat
the run as a blocking factor), at which point the slope equals the within-run correlation. To pool these
per-run correlations we use a random-effects (DerSimonian--Laird) meta-analysis of the per-run Fisher-$z$
values rather than an equal-weight mean, since runs vary in leaf count and noise. As an estimator-robustness check on a representative full-FT
condition (Llama-3.2-1B, translation probe, centroid DV), the pooling choices agree in sign and
magnitude---DL random-effects $r=-0.161$ ($95\%$ CI $[-0.28,-0.04]$, $p=0.011$), equal-weight mean
$r=-0.21$, one-stage mixed model with a random slope $r=-0.136$---with genuine between-run heterogeneity
($I^2\approx54\%$) that is not explained by leaf count. We report the DL random-effects estimate; the within-run partial
associations underlying it are shown per condition in Fig.~\ref{fig:regression_diagnostics}.


\section{Translation Probe Outputs for TaskHoldout Models}
\label{app:translation_probe}

The \texttt{runs\_no\_translation\_v2} condition trains models on summarization,
QA, and classification only, with translation withheld from the task pool. We use
the translation probe to test whether weight distance predicts behavioral similarity
even when the probe task was never trained. However, this condition reveals a
fundamental limitation of behavioral probing on out-of-distribution tasks: models
whose most recent fine-tuning was on classification have overwritten the base
model's translation ability entirely, producing degenerate single-token outputs
rather than translations.

Table~\ref{tab:translation_probe_examples} shows representative outputs for three
leaf nodes from a single run (\texttt{run\_001}). Nodes trained most recently on
summarization or QA retain the base model's translation ability and produce
coherent (if imperfect) French output. Nodes trained most recently on
classification output only the label vocabulary of their fine-tuning task,
regardless of input.

\begin{table}[h]
\caption{Translation probe outputs for three leaf nodes in \texttt{run\_001}
  (TaskHoldout condition, translation withheld from training).
  The same three English source sentences are shown for each node.
  Classification-trained nodes output label tokens only.}
\label{tab:translation_probe_examples}
\small
\setlength{\tabcolsep}{4pt}
\begin{center}
\begin{tabular}{p{1.8cm} p{1.6cm} p{8.0cm}}
\toprule
Node & Last task & Model output \\
\midrule
\multicolumn{3}{@{}p{11.6cm}@{}}{\textit{Source 1: ``And, most recently, it cannot excuse the failure to design a simple website more than three years since the Affordable Care Act was signed.''}} \\
n0  & QA/Summ.  & \textit{Et, plus r\'{e}cemment, il ne peut excuser le manque de conception d'une simple website plus de trois ann\'{e}es depuis le signe de la Loi sur le soins de sant\'{e}.} \\
n7  & Classification & \texttt{negative} \\
n8  & Classification & \texttt{entailment} \\
\midrule
\multicolumn{3}{@{}p{11.6cm}@{}}{\textit{Source 2: ``Since the Tribunal has declared it admissible, it would need to proceed to the second stage, the preliminary hearings.''}} \\
n0  & QA/Summ.  & \textit{Puisque le Tribunal a d\'{e}clar\'{e} admissible, il faudrait proc\'{e}der \`{a} la second \'{e}tape, les audiences pr\'{e}liminaires.} \\
n7  & Classification & \texttt{positive} \\
n8  & Classification & \texttt{entailment} \\
\midrule
\multicolumn{3}{@{}p{11.6cm}@{}}{\textit{Source 3: ``The bikes are jostling to get in front and take photos.''}} \\
n0  & QA/Summ.  & \textit{Les bikes sont jotling to get in front and take photos.} \\
n7  & Classification & \texttt{negative} \\
n8  & Classification & \texttt{entailment} \\
\bottomrule
\end{tabular}
\end{center}
\end{table}

Because classification-trained nodes produce degenerate outputs, pairwise semantic
similarity between nodes is driven by whether each model can translate at all,
not by genealogical proximity. This renders standard regression analysis
uninterpretable for this condition. We therefore treat the TaskHoldout condition
as a qualitative null: the absence of a meaningful weight--behavior link is
confirmed by construction, since behavioral similarity on the translation probe
is dominated by task identity (translating vs.\ not translating) rather than
lineage signal.

This degenerate-output collapse is specific to the encoder--decoder Flan-T5 TaskHoldout condition,
where classification-trained nodes overwrite translation entirely. It is \emph{not} a general property
of held-out probes. In the matched full-fine-tuning held-out conditions of
Table~\ref{tab:behavior_holdout} the nodes remain capable on the held-out task, and there the
weight$\to$behavior bridge is interpretable: negative and \emph{strengthening} for Llama-3.2-1B
(evidence of genuine lineage transmission), and only weakly positive for Flan-T5---consistent with the
residual task-identity confound described here rather than a clean reversal of the bridge.


\section{Behavioral Regression Diagnostics}
\label{app:regression_diagnostics}

Figure~\ref{fig:regression_diagnostics} shows the within-run partial association
between cosine weight distance and semantic output similarity for each fine-tuning
condition. Points are within-run demeaned (run fixed effects removed); the red line
is the OLS fit on the demeaned data. A negative slope indicates that model pairs
closer in weight space produce more similar translation outputs. The effect is
consistent in direction across all five Flan-T5 training conditions and the full-FT decoder Llama-3.2-1B
(centroid semantic DV; translation probe), with within-run $\bar r$ ranging from $-0.16$ to $-0.45$. The
held-out bridge, on the paired per-prompt DV and pooled by random-effects meta-analysis, is quantified
separately in \S\ref{sec:results} (Table~\ref{tab:behavior_holdout}).

\begin{figure}[h]
  \begin{center}
    \includegraphics[width=\linewidth]{\figpath/regression_diagnostics_semantic.png}
  \end{center}
  \caption{Within-run partial association between cosine weight distance and
    semantic output similarity (translation probe), for the five Flan-T5 fine-tuning
    conditions and the full-FT decoder Llama-3.2-1B (bottom right). Points are
    within-run demeaned; Fisher $z$-averaged within-run $r$ and run count shown per
    panel. All slopes are negative, confirming that weight proximity predicts
    behavioral similarity across training regimes and architectures.}
  \label{fig:regression_diagnostics}
\end{figure}

\end{document}
